% Part: first-order-logic
% Chapter: syntax-and-semantics
% Section: unique-readability

\documentclass[../../../include/open-logic-section]{subfiles}

\begin{document}

\olfileid{fol}{syn}{unq}

\olsection{唯一可读性}

\begin{explain}
公式的定义方式保证了每个公式都有\emph{唯一的读法}；
也就是说，根据公式的形成规则，本质上只有唯一一种构造它的方式，
也只有唯一一种“解释”它的方式。
若非如此，我们就会得到有歧义的公式，
即具有不止一种读法或解释的公式——这显然是我们想要避免的。
但更重要的是，若没有这个性质，我们将要给出的大多数定义和证明都无法进行。

或许阐明这一点最好的办法，是看看如果我们给出了不当的公式形成规则、
从而不能保证唯一可读性，会发生什么。
例如，我们可能忘了在联结词的形成规则中加括号，
比如，也许会允许这样：
\begin{quote}
如果 $!A$ 和 $!B$ 是公式，那么 $!A \lif !B$ 也是公式。
\end{quote}
从一个原子公式 $!D$ 出发，这将允许我们形成 $!D \lif !D$。
结合 $!D$，由此又可得到 $!D \lif !D \lif !D$。
但有两种方式可以做到这一点：
\begin{enumerate}
\item 令 $!A$ 为 $!D$，令 $!B$ 为 $!D \lif !D$。
\item 令 $!A$ 为 $!D \lif !D$，令 $!B$ 为~$!D$。
\end{enumerate}
相应地，公式~$!D \lif !D \lif !D$ 就有两种“读法”。
它既可看作形式为 $!B \lif !C$（其中 $!B$ 为 $!D$，$!C$ 为 $!D \lif !D$），
\emph{也可}看作形式为 $!B \lif !C$（其中 $!B$ 为 $!D \lif !D$，$!C$ 为~$!D$）的公式。

如果出现这种情况，我们的定义有时就会失效。
例如，在定义公式的主联结词时，我们会说：
在形如 $!B \lif !C$ 的公式中，主联结词是指明的那一个~$\lif$ 的出现。
但如果我们能用上述两种不同方式将公式 $!D \lif !D \lif !D$ 与 $!B \lif !C$ 匹配，
那么在第一种情况下，第一个 $\lif$ 的出现是主联结词，
而在第二种情况下，第二个 $\lif$ 的出现是主联结词。
然而，我们希望主联结词是公式的\emph{函数}，
即每个公式必须有且仅有一个主联结词出现。
\end{explain}

\begin{lem}
公式~$!A$ 中左括号和右括号的数量相等。
\end{lem}

\begin{proof}
我们通过对 $!A$ 的构造方式进行归纳来证明这一点。
这需要两步： 首先证明所有原子公式都具有所述性质（归纳基础）；
然后证明，在由已知公式构造新公式时，若已知公式满足该性质，则新公式也满足。

令 $l(!A)$ 为 $!A$ 中左括号的数量，$r(!A)$ 为 $!A$ 中右括号的数量；
$l(t)$ 与 $r(t)$ 分别表示项~$t$ 中左、右括号的数量。

\begin{prob}
    证明对任意项~$t$，$l(t) = r(t)$。
\end{prob}

\begin{enumerate}
\tagitem{prvFalse}{\indcase{!A}{\lfalse}{$\indfrm$ 的左括号数和右括号数均为 $0$。}}{}

\tagitem{prvTrue}{\indcase{!A}{\ltrue}{$\indfrm$ 的左括号数和右括号数均为 $0$。}}{}

\item \indcase{!A}{\Atom{R}{t_1,\dots,t_n}}{$l(\indfrm) = 1 + l(t_1) +
  \dots + l(t_n) = 1 + r(t_1) + \dots + r(t_n) = r(\indfrm)$。这里我们利用了这一事实（留作练习）：对任意项~$t$，$l(t) = r(t)$。}

\item \indcase{!A}{\eq[t_1][t_2]}{$l(\indfrm) = l(t_1) + l(t_2) =
  r(t_1) + r(t_2) = r(\indfrm)$。}

\tagitem{prvNot}{\indcase{!A}{\lnot !B}{由归纳假设，$l(!B) = r(!B)$。因此，$l(\indfrm) = l(!B) = r(!B) =
    r(\indfrm)$。}}{}

\item \indcase{!A}{(!B \ast !C)}{由归纳假设，$l(!B) =
  r(!B)$ 且 $l(!C) = r(!C)$。因此，$l(\indfrm) = 1 + l(!B) + l(!C) =
  1 + r(!B) + r(!C) = r(\indfrm)$。}

\tagitem{prvAll}{\indcase{!A}{\lforall[x][!B]}{由归纳假设，$l(!B) = r(!B)$。因此，$l(\indfrm) = l(!B) = r(!B) =
    r(\indfrm)$。}}{}

% Print case for \lexists if prvEx and not prvAll
\tagitem{prvAll,notprvEx}{}{\indcase{!A}{\lexists[x][!B]}{由归纳假设，$l(!B) = r(!B)$。因此，$l(\indfrm) = l(!B) = r(!B) =
    r(\indfrm)$。}}

% Just say similarly if prvEx and prvAll
\tagitem{notprvAll,notprvEx}{}{\indcase{!A}{\lexists[x][!B]}{类似可证。}}
\end{enumerate}
\end{proof}

\begin{defn}[真前缀]
若将符号串 $!B$ 与一个非空符号串拼接即可得到符号串~$!A$，则称 $!B$ 为 $!A$ 的\emph{真前缀}。
\end{defn}

\begin{lem}\ollabel{lem:no-prefix}
如果 $!A$ 是公式，且 $!B$ 是 $!A$ 的真前缀，那么 $!B$ 不是公式。
\end{lem}

\begin{proof}
练习。
\end{proof}

\begin{prob}
证明 \olref[fol][syn][unq]{lem:no-prefix}。
\end{prob}

\begin{prop}
\ollabel{prop:unique-atomic}
若 $!A$ 是原子公式，则它满足以下条件中的一个，且仅满足一个。
\begin{enumerate}
\tagitem{prvFalse}{$!A \ident \lfalse$.}{}
\tagitem{prvTrue}{$!A \ident \ltrue$.}{}
\item $!A \ident \Atom{R}{t_1,\dots,t_n}$ 其中 $R$ 为 $n$ 元谓词符号，$t_1$，\dots，$t_n$ 为项，且 $R$、$t_1$、\dots、$t_n$ 均被唯一确定。
\item $!A \ident \eq[t_1][t_2]$ 其中 $t_1$ 和 $t_2$ 为唯一确定的项。
\end{enumerate}
\end{prop}

\begin{proof}
（留作）练习。
\end{proof}

\begin{prob}
证明 \olref[fol][syn][unq]{prop:unique-atomic}（提示：针对项来表述并证明 \olref[fol][syn][unq]{lem:no-prefix} 的一个对应版本。）
\end{prob}

\begin{prop}[唯一可读性]
每个公式都满足以下条件中的一个，且仅满足一个。
\begin{enumerate}
\item $!A$ 是原子公式。

\tagitem{prvNot}{$!A$ 具有形式 $\lnot !B$.}{}

\tagitem{prvAnd}{$!A$ 具有形式 $(!B \land !C)$.}{}

\tagitem{prvOr}{$!A$ 具有形式 $(!B \lor !C)$.}{}

\tagitem{prvIf}{$!A$ 具有形式 $(!B \lif !C)$.}{}

\tagitem{prvIff}{$!A$ 具有形式 $(!B \liff !C)$.}{}

\tagitem{prvAll}{$!A$ 具有形式 $\lforall[x][!B]$.}{}

\tagitem{prvEx}{$!A$ 具有形式 $\lexists[x][!B]$.}{}
\end{enumerate}
此外，在每种情形下，$!B$（或 $!B$ 与 $!C$）均是唯一确定的。
这意味着，例如，不存在两组不同的 $!B, !C$ 与 $!B', !C'$，使得 $!A$ 同时具有形式
\iftag{prvIf}{$(!B \lif !C)$ 和 $(!B' \lif !C')$.}{
    \iftag{prvOr}{$(!B \lor !C)$ 和 $(!B' \lor !C')$.}{
      \iftag{prvAnd}{$(!B \land !C)$ 和 $(!B' \land !C')$.}{}}}
\end{prop}

\begin{proof}
形成规则要求：如果公式不是原子的，则它必须以左括号 (、\iftag{prvNot}{$\lnot$,}{}或量词开头。
另一方面，以谓词符号、函数符号、常项符号\iftag{prvFalse}{, $\lfalse$}{}\iftag{prvTrue}{, $\ltrue$}{}之一开头的公式必为原子公式。

因此，我们实际上只需证明：如果 $!A$ 具有形式 $(!B \ast !C)$ 并且也具有形式 $(!B' \mathbin{\ast'} !C')$，那么 $!B \ident !B'$、$!C \ident !C'$，且 $\ast = {\ast'}$。

假设 $!A \ident (!B \ast !C)$ 且 $!A \ident (!B'
\mathbin{\ast'} !C')$ 同时成立。
那么要么 $!B \ident !B'$，要么 $!B \not\ident !B'$。
若是前者，显然 $\ast = {\ast'}$ 且 $!C \ident !C'$，因为此时 $!B$ 与 $!B'$ 是 $!A$ 中从同一位置开始且长度相同的子串。
另一种情况是 $!B \not\ident !B'$。
既然 $!B$ 和 $!B'$ 同为 $!A$ 中起始于同一位置的子串，则其中一个必为另一个的真前缀。
但由 \olref{lem:no-prefix} 可知，这是不可能的。
\end{proof}

\end{document}